% Lösungen Einheit 1 von 3 – Potenz-, Faktor- und Summenregel (zusammenbau v0.1)
\einheitenkopf{Einheit 1 von 3 · Potenz-, Faktor- und Summenregel}

\erg{10}{a) $a$ – der Parameter bleibt beim Ableiten stehen wie eine Zahl.}
\erg{11}{a) Potenzregel – es gibt keinen Vorfaktor und keine Summe. \quad b) $f'(x) = 6x^2 + 10x - 4$ \quad c) $f'(x) = \frac{1}{2}x^2 + 3x - 2$ \quad d) $f'(x) = -10x^4 + 13x^3 + 18x^2 - 60x$; $f''(x) = -40x^3 + 39x^2 + 36x - 60$; $f'''(x) = -120x^2 + 78x + 36$ \quad e) $f'(x) = -15x^4 + 12x^2 - 1$; $F(x) = -\frac{1}{2}x^6 + x^4 - \frac{1}{2}x^2$ \quad f) $f'(x) = x^2 - \frac{9}{2}$; $f''(x) = 2x$; $f'''(x) = 2$; Anstieg: $f'(0) = -\frac{9}{2} = -4{,}5$ \quad g) $f_a'(x) = 4x^3 - 3a x^2 = x^2 \cdot (4x - 3a)$ \quad h) $f'(x) = 3x^2 - 18x + 24$; ausmultipliziert: $3 \cdot (x - 2) \cdot (x - 4) = 3 \cdot (x^2 - 6x + 8) = 3x^2 - 18x + 24$ – beide stimmen überein. \quad i) $f'(x) = x^4 - 18x^2 + 81 = (x^2 - 9)^2 = ((x - 3) \cdot (x + 3))^2 = (x - 3)^2 \cdot (x + 3)^2$ \quad j) $f'(x) = x^3 + 3x^2 - 4x$; $f''(x) = 3x^2 + 6x - 4$; $f'''(x) = 6x + 6$; $f(2) = 4$, also liegt $P$ auf dem Graphen; $f'(2) = 12 \neq 0$, also ist $P$ kein Extrempunkt.}
\erg{12}{a) für $x \to +\infty$: $f(x) \to +\infty$; für $x \to -\infty$: $f(x) \to -\infty$; $f'(x) = 6x^2 + 2x - 12$; $f''(x) = 12x + 2$; $f'''(x) = 12$}
\erg{13}{a) Die Exponenten wurden nicht um eins gesenkt; richtig: $f'(x) = 12x^3 - 4x$. \quad b) Der Graph ist eine waagerechte Gerade; ihre Steigung ist an jeder Stelle null.}
